\documentclass[11pt]{article}
\usepackage[utf8]{inputenc}
\usepackage[T1]{fontenc}
\usepackage{lmodern}
\usepackage[spanish]{babel}
\usepackage{amsmath,amssymb}
\usepackage{graphicx}
\usepackage[left=2cm,right=2cm,top=2cm,bottom=2cm]{geometry}
\usepackage{enumitem}
\usepackage{booktabs}
\usepackage{array}
\usepackage{xcolor}
\usepackage{framed}

\spanishdecimal{.}

\setlist[enumerate,1]{label=\textbf{\arabic*.}, leftmargin=*, itemsep=6pt}
\setlist[enumerate,2]{label=(\alph*), itemsep=4pt}
\setlength{\parindent}{0pt}
\setlength{\parskip}{4pt}
\pagestyle{plain}

\newcommand{\Ein}{E_{\text{in}}}
\newcommand{\Eout}{E_{\text{out}}}
\newcommand{\Hyp}{\mathcal{H}}
\newcommand{\dvc}{d_{\text{VC}}}
\newcommand{\vx}{\mathbf{x}}
\newcommand{\vw}{\mathbf{w}}
\newcommand{\vy}{\mathbf{y}}
\newcommand{\R}{\mathbb{R}}
\newcommand{\T}{^{\top}}
\newcommand{\E}{\mathbb{E}}
\newcommand{\Prob}{\mathbb{P}}
\newcommand{\pts}[1]{\hfill\textbf{(#1~puntos)}}
\newcommand{\linea}[1]{\rule{#1}{0.4pt}}
\newcommand{\ptsinc}[1]{{\unskip\nobreak\hfil\penalty50\hskip1em\hbox{}\nobreak\hfill(#1)\parfillskip=0pt\finalhyphendemerits=0\par}}

\newcommand{\encabezado}[2]{%
\begin{center}
{\large\bfseries Reconocimiento de Patrones}\\[2pt]
Universidad de Sonora \textperiodcentered{} Licenciatura en Ciencias de la Computación\\
Semestre 2027-1 \textperiodcentered{} Prof.\ Julio Waissman Vilanova\\[10pt]
{\Large\bfseries #1}%
\ifdefined\tipodoc\\[4pt]{\large\itshape\tipodoc}\fi
\end{center}
\ifdefined\tipodoc\else
\vspace{6pt}
\textbf{Nombre:} \hrulefill\quad \textbf{Expediente:} \rule{3.5cm}{0.4pt}\par
\vspace{6pt}
\fi
#2
\par\medskip\hrule\medskip
}

\newcommand{\espacio}[1]{\par\vspace{#1}}
\newcommand{\nuevapagina}{\newpage}


\begin{document}

\encabezado{Examen rápido 2 \textperiodcentered{} Aprendizaje no supervisado}{%
Tiempo: 20 minutos. Responde en esta hoja. Puedes usar calculadora; no se permiten apuntes. Deja indicadas las operaciones.}

\begin{enumerate}

%% ---------------------------------------------------------------
\item \textbf{k-medias a mano.} Tenemos seis puntos en el plano:
\[
\vx_1 = (1, 1),\quad \vx_2 = (2, 1),\quad \vx_3 = (1, 3),\quad \vx_4 = (5, 4),\quad \vx_5 = (6, 5),\quad \vx_6 = (7, 3).
\]
Corremos el algoritmo de Lloyd con $k = 2$ e iniciamos los centros en $\mu_1 = \vx_1$ y $\mu_2 = \vx_3$. Sea $J = \sum_i \lVert \vx_i - \mu_{c(i)} \rVert^2$.

\begin{enumerate}
  \item Haz el primer paso de asignación. ¿Qué puntos quedan con cada centro y cuánto vale $J$ con esa asignación y los centros iniciales?
  \espacio{3.2cm}

  \item Actualiza los centros y repite asignación y actualización hasta que nada cambie. Escribe los centros finales, la asignación final y el valor final de $J$.
  \espacio{5.5cm}

  \item ¿Por qué podemos estar seguros de que $J$ nunca aumenta de un paso al siguiente? ¿Garantiza eso que encontramos el agrupamiento con el menor $J$ posible?
  \espacio{2.5cm}
\end{enumerate}

\nuevapagina
%% ---------------------------------------------------------------
\item \textbf{¿A qué componente pertenece?} Una mezcla de dos gaussianas en una dimensión tiene componentes $\mathcal N(0, 1)$ y $\mathcal N(4, 1)$ (medias 0 y 4, ambas con varianza 1). Recuerda que
\[
\gamma_{2}(x) = \frac{\pi_2\, \mathcal N(x \mid \mu_2, \sigma^2)}{\pi_1\, \mathcal N(x \mid \mu_1, \sigma^2) + \pi_2\, \mathcal N(x \mid \mu_2, \sigma^2)} .
\]
\begin{enumerate}
  \item Con $\pi_1 = \pi_2 = 0.5$, calcula la responsabilidad de la segunda componente para $x = 2.5$. \emph{Sugerencia:} escribe $\gamma_2 = 1/(1 + e^{-a})$ y calcula $a$.
  \espacio{3cm}

  \item Repite con $\pi_1 = 0.8$, $\pi_2 = 0.2$.
  \espacio{2cm}

  \item En cada caso, ¿para qué $x$ vale la responsabilidad exactamente $0.5$? ¿Hacia dónde se mueve esa frontera al cambiar las proporciones y por qué tiene sentido?
  \espacio{3cm}
\end{enumerate}

%% ---------------------------------------------------------------
\item \textbf{PCA con una matriz de $2 \times 2$.} Los datos centrados de dos atributos tienen matriz de covarianza
\[
S = \begin{pmatrix} 5 & 2 \\ 2 & 2 \end{pmatrix}.
\]
\begin{enumerate}
  \item Encuentra los valores propios de $S$ y la primera componente principal $\mathbf u_1$ (unitaria). ¿Qué fracción de la varianza explica?
  \espacio{3.5cm}

  \item Para el punto centrado $\vx = (3, 1)$, calcula su coordenada $z_1$ en la primera componente, su reconstrucción $\hat\vx = z_1 \mathbf u_1$ y el error $\lVert \vx - \hat\vx \rVert^2$.
  \espacio{3cm}
\end{enumerate}

\end{enumerate}


\end{document}
